% =====================================================================
% SECCION 6 -- DISENO DEL SISTEMA
% Responsable: no corresponde al frente etico-legal.
%
% NOTA PARA EL EQUIPO: las decisiones tecnicas que se documenten aqui son
% auditadas en la Seccion \ref{sec:auditoria} del marco etico-legal.
% Cualquier cambio de arquitectura, de dataset o de politica de retencion
% debe comunicarse al frente etico-legal, porque cambia la auditoria.
% =====================================================================
% 14/09/2026: reestructuracion al indice del modelo del curso. Recibe el
% apartado "Arquitectura propuesta" de la antigua Metodologia; el resto de
% aquella seccion esta en secciones/desarrollo_modelo.tex (Seccion 8).
%
% 24/09/2026: la seccion deja de ser esqueleto.
%   6.1 Los seis componentes de procesamiento se trasladan en prosa desde
%       la Seccion 6.1 de INFORME_COMPLETO_ANGEL.md (misma sustancia; se
%       quitan las referencias a casos de uso y a puntos de PENDIENTES.md).
%   6.2 a 6.4 Componentes nuevos de la consola, documentados con las
%       capturas de imagenes/aerotrack-01 ... aerotrack-17.
% Cada figura describe solo lo que se ve en la captura. Las de los pasos
% 4, 6 y 8 del asistente salen sin video porque aun no hay grabaciones del
% terminal; el texto lo dice en una sola frase (apartado 6.4).
%
% \figuraconsola{archivo}{pie}{etiqueta}
% Todas las capturas van centradas y al mismo ancho (0,85 del ancho de
% texto). Para cambiar el ancho de todas a la vez, se cambia aqui.
% [H] (paquete float, cargado en main.tex) fija cada figura tras su
% parrafo, para que el recorrido paso a paso se lea en orden.
\newcommand{\figuraconsola}[3]{%
  \begin{figure}[H]
    \centering
    \includegraphics[width=0.85\textwidth]{imagenes/#1}
    \caption{#2}
    \label{#3}
  \end{figure}}

\section{Diseño del Sistema}
\label{sec:diseno}
\label{sec:metodologia}

Esta sección describe el diseño del prototipo AeroTrack, la consola con
la que se configura y opera el sistema de estimación de aglomeraciones.
El apartado~\ref{sec:arquitectura} presenta los seis componentes que
organizan el procesamiento de las imágenes. Los apartados siguientes
documentan los componentes incorporados en la versión actual de la
consola: la gestión de usuarios por roles, la gestión de proyectos
independientes, el asistente de configuración en ocho pasos y las alertas
de seguridad con aviso por correo.

\subsection{Arquitectura propuesta}
\label{sec:arquitectura}

La Sección~\ref{sec:sintesis-antecedentes} deja planteada la elección
entre los tres enfoques revisados y, sobre todo, la granularidad de la
salida, que es la decisión con mayor consecuencia para el análisis de
proporcionalidad. El prototipo organiza el procesamiento en seis
componentes, cada uno con una responsabilidad única.

El primero es la adquisición, que lee cada fuente de cámara, sea un
archivo de video, una cámara USB o un flujo RTSP. El segundo es la
detección: un detector intercambiable (HOG, P2PNet o YOLO) produce
posiciones en píxeles con una confianza asociada. La elección del
detector la hace el operador de forma explícita y no el sistema, porque
ninguno de los tres es adecuado para todos los casos. El tercero es el
seguimiento dentro de cada cámara. Con los detectores HOG y P2PNet es una
adaptación de ByteTrack~\cite{Zhang_2022_ECCV} a puntos, con filtro de
Kalman de velocidad constante y asignación húngara en dos etapas según la
confianza; con YOLO se sigue la caja completa de cada persona.

El cuarto componente es la calibración y proyección: cada cámara tiene
una homografía propia hacia el plano, que se valida por consistencia
geométrica antes de aceptarse. El quinto es la fusión entre cámaras, que
asocia identidades por proximidad espacio-temporal en el plano compartido
y solo recurre a un descriptor de apariencia no biométrico como desempate
de última instancia. El sexto reúne persistencia, analítica y
presentación: mantiene en memoria el estado agregado de la sesión y lo
expone en la consola, que ofrece configuración, plano en vivo, video por
cámara y reportes. Los apartados siguientes describen cómo se presenta
esa consola al usuario.

\subsection{Gestión de usuarios y roles}
\label{sec:diseno-usuarios}

El acceso a la consola requiere usuario y contraseña. En la primera
ejecución el sistema pide crear un primer usuario
(Figura~\ref{fig:primer-usuario}), que queda con el rol de operador y
puede dar de alta a los demás. La contraseña debe tener al menos seis
caracteres y, según indica la propia pantalla, se guarda cifrada en el
equipo y nunca en texto plano.

\figuraconsola{aerotrack-01-primer-usuario.png}
  {Creación del primer usuario en la primera ejecución de la consola.}
  {fig:primer-usuario}

Una vez dentro, la vista general muestra en la cabecera el proyecto
activo junto al usuario y su rol, y organiza la consola en cinco
pestañas: Vista general, Alertas, Videos y resultados, Reportes y
Configuración (Figura~\ref{fig:vista-general}). Desde esta vista se
filtra por nivel, cámara, zona de imagen y rendimiento, y se inicia el
monitoreo sobre el plano del nivel. En la captura los contadores de
cámaras seleccionadas, entradas, salidas y tiempo de análisis están en
cero porque el proyecto aún no tiene cámaras configuradas.

\figuraconsola{aerotrack-02-vista-general.png}
  {Vista general de la consola, con el proyecto activo en la cabecera y el
   plano del nivel 3.}
  {fig:vista-general}

El sistema distingue dos roles (Figura~\ref{fig:usuarios}). El operador
configura cámaras y plano; el administrador consulta resultados y ajusta
los umbrales de alerta, sin acceso a la configuración. El panel de
usuarios lista las cuentas existentes con su rol, permite eliminarlas y
permite crear otras indicando nombre, contraseña y rol.

\figuraconsola{aerotrack-03-usuarios.png}
  {Panel de usuarios del sistema y formulario de alta con selección de rol.}
  {fig:usuarios}

\subsection{Gestión de proyectos}
\label{sec:diseno-proyectos}

Un proyecto representa un espacio monitoreado, por ejemplo un nivel de un
terminal. Cada proyecto guarda por separado su propio plano, sus cámaras
y su calibración, de modo que la configuración de un espacio no se mezcla
con la de otro y es posible mantener varios proyectos y cambiar de uno a
otro cuando haga falta. Al entrar en la configuración, la consola ofrece
dos acciones: crear un proyecto nuevo, que conduce al asistente descrito
en el apartado~\ref{sec:diseno-asistente}, o añadir una cámara al
proyecto abierto (Figura~\ref{fig:crear-proyecto}). Para crear el
proyecto basta con darle un nombre.

\figuraconsola{aerotrack-04-crear-proyecto.png}
  {Opciones para crear un proyecto nuevo o añadir una cámara al proyecto
   abierto.}
  {fig:crear-proyecto}

Los proyectos guardados se listan con su nombre, el espacio al que
corresponden, el número de cámaras y de planos, el momento de la última
edición y su estado (Figura~\ref{fig:proyectos-guardados}). En la captura
figuran dos proyectos, ambos marcados como incompletos: el proyecto
principal, sobre el nivel 3 del aeropuerto y con cuatro planos, y el
proyecto Terminal~A, Nivel~1, con un plano. Desde la lista se abre un
proyecto, se retoma su configuración en el punto en que quedó, se
renombra directamente sobre su tarjeta (Figura~\ref{fig:renombrar}) o se
elimina.

\figuraconsola{aerotrack-05-proyectos-guardados.png}
  {Lista de proyectos guardados, cada uno con sus cámaras, planos y estado.}
  {fig:proyectos-guardados}

\figuraconsola{aerotrack-06-renombrar-proyecto.png}
  {Cambio de nombre de un proyecto sobre su propia tarjeta.}
  {fig:renombrar}

Al entrar en la consola, una pantalla pregunta en qué proyecto se quiere
trabajar, muestra para cada uno su espacio, su número de cámaras y el
estado de su configuración, y señala el último abierto
(Figura~\ref{fig:selector-proyecto}). Así el usuario elige de forma
expresa el espacio en el que va a trabajar y no mezcla datos de dos
lugares distintos.

\figuraconsola{aerotrack-07-selector-proyecto.png}
  {Selección del proyecto de trabajo al entrar en la consola.}
  {fig:selector-proyecto}

\subsection{Asistente de configuración}
\label{sec:diseno-asistente}

La configuración de un proyecto se completa con un asistente de ocho
pasos: proyecto, plano, conexión, prueba, ubicación, calibración, zonas y
alertas, y revisión. Una barra superior muestra los pasos, marca los
completados y lleva la cuenta de cuántos están listos, y una indicación
en la cabecera confirma que el avance queda guardado, de modo que un
proyecto incompleto puede retomarse desde la lista de proyectos. Las
capturas de los pasos 4, 6 y 8 se tomaron sin video, porque en el estado
actual del prototipo aún no se dispone de grabaciones del terminal.

\subsubsection*{Paso 1: proyecto}

Se identifica el espacio con el nombre del aeropuerto o proyecto y el del
terminal y nivel, y se elige el tipo de fuente: videos grabados, para
analizar grabaciones y consultar resultados por periodo, o cámaras en
vivo, conectadas por USB, RTSP o IP (Figura~\ref{fig:paso1}). Una nota
lateral recuerda que el análisis es de presencia y recorridos, sin
reconocimiento facial.

\figuraconsola{aerotrack-08-paso1-proyecto.png}
  {Paso 1 del asistente: datos del espacio y tipo de fuente.}
  {fig:paso1}

\subsubsection*{Paso 2: plano}

Se define el plano del espacio: se elige uno de los planos del
aeropuerto, se importa una imagen o un PDF, o se dibuja el espacio con
rectángulos, polígonos y líneas. El plano se acompaña de su ancho y su
alto en metros, y admite marcar un área de trabajo como referencia
opcional (Figura~\ref{fig:paso2}).

\figuraconsola{aerotrack-09-paso2-plano.png}
  {Paso 2 del asistente: plano del nivel y sus dimensiones en metros.}
  {fig:paso2}

\subsubsection*{Paso 3: conexión}

Se añaden las cámaras del proyecto (Figura~\ref{fig:paso3}). Para cada
una se indica el plano asociado, el nombre visible en el mapa, la
ubicación, la orientación y la fuente, con la ruta del video o la
dirección de la cámara, y se fijan el umbral de alerta en personas y el
tiempo de concentración sostenida (Figura~\ref{fig:paso3-camara}). En el
mismo paso se marca la zona útil de la imagen, que deja fuera reflejos,
otras plantas y áreas que no deben analizarse, y, si la cámara ve una
puerta, un acceso formado por un área exterior y otra interior, entre las
que se cuenta el paso de personas y no visitantes únicos
(Figura~\ref{fig:paso3-zona}).

\figuraconsola{aerotrack-10-paso3-conexion.png}
  {Paso 3 del asistente: lista de cámaras del proyecto, todavía vacía.}
  {fig:paso3}

\figuraconsola{aerotrack-11-paso3-camara.png}
  {Paso 3 del asistente: datos de la cámara, fuente y umbrales de alerta.}
  {fig:paso3-camara}

\figuraconsola{aerotrack-12-paso3-zona-util.png}
  {Paso 3 del asistente: zona útil de la imagen y acceso de un local.}
  {fig:paso3-zona}

\subsubsection*{Paso 4: prueba}

Se comprueba la cámara en movimiento para revisar encuadre, luz y
retardo. El panel muestra el estado de la fuente, su resolución, su
frecuencia de cuadros y si está calibrada, y permite iniciar una prueba
de seguimiento que presenta el video procesado con identificadores
temporales (Figura~\ref{fig:paso4}).

\figuraconsola{aerotrack-13-paso4-prueba.png}
  {Paso 4 del asistente: prueba de la cámara, con el visor sin señal.}
  {fig:paso4}

\subsubsection*{Paso 5: ubicación}

Se coloca la cámara sobre el plano en el punto donde está instalada y se
orienta hacia donde apunta, con su dirección, alcance, altura y ancho de
cobertura (Figura~\ref{fig:paso5}). Esa cobertura es orientativa: la
precisión la aporta la calibración del paso siguiente.

\figuraconsola{aerotrack-14-paso5-ubicacion.png}
  {Paso 5 del asistente: posición y cobertura de la cámara sobre el plano.}
  {fig:paso5}

\subsubsection*{Paso 6: calibración}

Se relaciona el suelo del video con el plano marcando puntos de
referencia, como esquinas o cruces de baldosas, primero en la imagen y
después en el mapa (Figura~\ref{fig:paso6}). Se requieren al menos cuatro
puntos repartidos; con cinco o más, cada referencia se comprueba además
dejándola fuera del ajuste.

\figuraconsola{aerotrack-15-paso6-calibracion.png}
  {Paso 6 del asistente: calibración del suelo del video con el plano.}
  {fig:paso6}

\subsubsection*{Paso 7: zonas y alertas}

Se dibujan sobre el plano las zonas de interés con polígonos, rectángulos
o círculos, y se define cuándo avisar de una aglomeración mediante un
radio en metros, un número mínimo de personas y un tiempo de permanencia
(Figura~\ref{fig:paso7}). La pantalla advierte que una fila y una
aglomeración pueden parecerse geométricamente, por lo que los límites
deben ajustarse al uso de cada zona.

\figuraconsola{aerotrack-16-paso7-zonas.png}
  {Paso 7 del asistente: trazado de una zona y parámetros de aviso de
   aglomeración.}
  {fig:paso7}

\subsubsection*{Paso 8: revisión}

Se revisa el proyecto antes de abrir el monitoreo: una lista indica qué
está preparado y qué sigue pendiente, y un resumen muestra el número de
cámaras, zonas dibujadas y fuentes probadas, junto con la escala
(Figura~\ref{fig:paso8}). En la captura, con la cámara sin probar ni
calibrar, la opción de guardar y abrir el monitoreo aparece desactivada y
queda disponible la de continuar las pruebas en laboratorio. La pantalla
aclara que una verificación completada no sustituye la validación de
precisión.

\figuraconsola{aerotrack-17-paso8-revision.png}
  {Paso 8 del asistente: revisión del proyecto antes de abrir el monitoreo.}
  {fig:paso8}

\subsection{Alertas de seguridad con aviso por correo}
\label{sec:diseno-alertas}

La vista de alertas reúne lo que necesita el personal de seguridad para
atender una aglomeración (Figura~\ref{fig:alertas-correo}). El apartado
\enquote{Cuándo avisar} fija los umbrales que definen una aglomeración: el
número mínimo de personas, la permanencia en segundos y el radio en
metros; en la captura, 4 personas, 3 segundos y 1,5 metros. Estos
umbrales se cambian desde la misma vista, sin entrar en la configuración,
y son el único ajuste que también puede hacer el administrador.

El panel \enquote{Avisar por correo}, desplegado en la captura y todavía
desactivado, pide la cuenta que envía, su contraseña de aplicación y los
correos que reciben el aviso, separados por coma, y permite enviar un
correo de prueba. Una nota aclara que el servicio de correo no acepta la
contraseña normal de la cuenta, sino una contraseña de aplicación
generada con la verificación en dos pasos. Debajo, tres contadores
muestran las alertas activas, el total de la sesión y los incidentes sin
revisar, y la bitácora de incidentes conserva cada caso aunque la sesión
termine, para que el personal lo marque cuando alguien haya ido a
revisarlo. En la captura la bitácora aparece sin registros porque en esta
fase no se han procesado grabaciones del terminal.

Con el aviso activado, el sistema envía un correo por cada incidente
nuevo de aglomeración o de equipaje sin custodia, con el lugar, la cámara
y el motivo del aviso, y la advertencia de que requiere comprobación
humana; un aviso ya enviado no se reenvía. La contraseña se guarda solo
en el equipo, fuera de los archivos del proyecto y del control de
versiones: la interfaz nunca la muestra de vuelta y no se copia al
duplicar un proyecto.

\figuraconsola{aerotrack-18-alertas-correo.png}
  {Vista de alertas: umbrales de aglomeración, aviso por correo desplegado,
   contadores y bitácora de incidentes.}
  {fig:alertas-correo}
